% ===== §14 Conclusions =====
\section{Conclusions: At the Boundaries of the Frameworks}
\label{sec:conclusions}

\noindent
This paper has travelled a considerable distance from its starting point: a flower on a path, shared with the person one loves, and the happiness that arises from that sharing. The distance is not a departure from the starting point but a return to it---a return in which the flower is still the flower it was, the sharing is still the sharing it was, but both are now seen in the light of a theoretical understanding that was not available at the beginning. This is the spiral structure that \pinkref{Paper~IX} identified as the form of relational development and that the series, in its own form, enacts \citep{paperix}: the return to the beginning that carries back the holonomy of everything traversed between departure and return.

In keeping with the series' commitment to honest multi-framework analysis---and with the formal律令established in \pinkref{Paper~IX} that genuine conclusions are polyphonic, that each framework says what it can say and honestly acknowledges what it cannot---this concluding section does not offer a synthetic summary that collapses the paper's arguments into a single unified thesis. Instead, it presents the conclusions that each of the paper's major frameworks can and cannot reach, and ends with a reflection on what the paper's incompleteness reveals about the nature of the phenomenon it has been studying.

\subsection{What the Philosophical Framework Can and Cannot Say}
\label{subsec:conc_phil}

\noindent
The philosophical framework of the GRB account---its relational ontology, its account of the \emph{between} as a positive ontological region, its reconception of contingency as belonging primarily to the relation rather than to the individual---can say the following with confidence:

The primary subject of relational happiness is the relational field, not the individual. The happiness that arises from sharing the flower on the path is not the sum of two individual happinesses but an emergent property of the relational field: a property that belongs to the \emph{between} and that cannot be located in either participant. The act of sharing is not the communication of a private experience but the constitutive act by which a contingent event enters the relational field and is transformed into a relational event. The subject of contingency is the relation, not the individual.

The philosophical framework cannot say, from its own resources alone, why the relational field has the capacity to generate this emergent happiness---what the mechanism is, how the co-evolutionary dynamics work at the level of neural implementation, what the formal conditions for the emergence of shared attractors are. These questions require the formal and empirical frameworks that the paper has also developed. The philosophical framework establishes the conceptual space within which the formal and empirical questions can be asked; it does not answer them.

The philosophical framework also cannot resolve the tension between the GRB account's relational ontology and the genuine phenomenological reality of individual experience. The individual does experience their own happiness; there is something that it is like to be this particular person noticing the flower and feeling the happiness of sharing it. The GRB account holds that this individual experience is the relational field's happiness perceived from one of its constituted perspectives, but this claim, however philosophically well-grounded, does not fully dissolve the tension between the primacy of the relational field and the undeniable reality of individual experience. The tension is real, and its resolution---if it has a resolution---lies beyond what the philosophical framework of the present paper can provide.

\subsection{What the Formal Framework Can and Cannot Say}
\label{subsec:conc_formal}

\noindent
The formal dynamical systems framework---the coupled oscillator model, the Kuramoto synchronisation analysis, the quantum entanglement structural analogy, the relational STDP mechanism, the phase space analysis of shared attractors---can say the following:

The co-evolutionary dynamics of an intimate relational system are real, formally well-defined, and distinct from the sum of the individual dynamics that compose them. The shared attractors of a coupled dynamical system are not the attractors of either individual system; they are genuine emergent properties of the coupling. The relational STDP mechanism provides a precise account of how contingent events produce structural modifications of the relational field, and it implies specific, testable predictions about the temporal dynamics of sharing and the cumulative trajectory of relational development.

The formal framework cannot say what the felt quality of the shared happiness is---cannot, from its equations and phase diagrams, generate the phenomenological description of \xref{sec:phenomenology}. The formal framework operates at the level of dynamics and structure; it does not operate at the level of experience. The relationship between the formal level and the phenomenological level---the hard problem, in its relational form---remains open, and the present paper does not solve it. It is enough to have shown that the formal and phenomenological levels are complementary descriptions of the same phenomenon, without claiming that either reduces to the other.

The formal framework also cannot guarantee that the structural isomorphisms it has identified---between the relational system and coupled oscillators, between the relational system and quantum entanglement, between the relational system and spiking neural networks---are more than heuristically useful analogies. The formal framework has been presented throughout as providing structural isomorphisms rather than literal identifications; the limits of this claim must be honestly acknowledged. The relational system is not a quantum system, not a spiking neural network, not a Kuramoto oscillator. It is a system that shares certain formal properties with these well-studied systems, and the shared formal properties illuminate aspects of the relational system's dynamics that would otherwise be difficult to describe with precision. But the illumination is partial, and the limits of the analogies are as important as their applicability.

\subsection{What the Empirical Framework Can and Cannot Say}
\label{subsec:conc_empirical}

\noindent
The empirical framework---the hyperscanning methodology, the longitudinal dyadic research programme, the physiological synchrony measurements, the behavioural coding protocols---can say the following:

There are empirically testable predictions that follow from the GRB account, and these predictions distinguish it from the individualist alternatives. The prediction that shared contingent events produce greater inter-brain synchrony than the same events experienced independently, the prediction that timely sharing produces stronger physiological coupling than delayed sharing, the prediction that longitudinal co-evolutionary dynamics produce cumulative structural modifications of the relational attractor landscape: these are all empirically tractable claims that can be tested with existing or developable methodology.

The empirical framework cannot confirm the philosophical and phenomenological claims of the GRB account---cannot show, from EEG data or fMRI activations, that the primary subject of happiness is the relational field rather than the individual. The empirical framework operates at the level of neural and physiological correlates; it does not operate at the level of ontological claims. The relationship between the empirical evidence and the philosophical framework is one of mutual constraint rather than unidirectional grounding: the empirical findings constrain the philosophical account by ruling out accounts that are inconsistent with the data, but they do not uniquely determine the philosophical account by entailing it.

The empirical framework also cannot address the cultural and cross-cultural dimensions of the phenomenon without extending its methodology beyond the WEIRD samples that currently dominate the hyperscanning literature. The cross-cultural findings reported in \xref{sec:culture} are philosophical and anthropological rather than empirical; a genuinely adequate empirical research programme would need to test the GRB account's predictions across the cultural contexts identified in that section, and would need to develop methodology sensitive to the culturally specific forms that relational co-evolution takes in different traditions.

\subsection{What the Hermeneutic Re-reading Can and Cannot Say}
\label{subsec:conc_herm}

\noindent
The hermeneutic re-reading of the classical happiness concepts---the GRB transformations of virtue, \emph{phronesis}, \emph{philia}, flow, \emph{ataraxia}, \emph{apatheia}, Maslow's hierarchy, SDT, Frankl's meaning, subjective well-being, and PERMA---can say the following:

Each of the classical concepts is transformed, rather than simply refuted, by the relocation of the subject of flourishing from the individual to the relational field. The classical concepts retain their genuine insights---the Aristotelian account of virtue as developed through practice, the Stoic insight into the relationship between vulnerability and suffering, Csikszentmihalyi's account of the skill-challenge balance as a condition for optimal experience---while shedding the individualist presuppositions that have limited their application to shared happiness. The hermeneutic result is not the rejection of the classical tradition but its extension and completion.

The hermeneutic re-reading cannot claim that the GRB transformations of the classical concepts are the only possible transformations, or that they are fully adequate to the richness of the original concepts. Each subsection of \xref{sec:hermeneutics} has been, of necessity, a compressed and selective reading of concepts that have generated centuries of philosophical commentary; the GRB transformation has focused on the single dimension most relevant to the present paper's argument, and has inevitably missed dimensions of the original concepts that a more complete treatment would address. The hermeneutic task this paper has begun is not complete; it is an opening of a conversation that future work must continue.

\subsection{Three Concentric Circles}
\label{subsec:circles}

\noindent
The paper closes its conclusions with an image that the framework does not generate but that the framework's convergent findings suggest. The happiness of the shared flower, examined from the multiple perspectives of this paper, appears in the form of three concentric circles.

The innermost circle is the moment itself: the contingent event, the act of sharing, the resonance of joint attention, the felt quality of the happiness that arises from the \emph{between}. This circle is the phenomenological circle---the circle of experience, of the living moment in which the relational field's co-evolutionary activity presents itself to the individuals who participate in it. It is the circle that the series' literary envois have always inhabited, from the rose of \pinkref{Paper~IX} to the red silk of \pinkref{Paper~XIII} \citep{paperix,paperxiii}.

The middle circle is the relational field: the accumulated history of shared contingent events, the shared attractor landscape built through a lifetime of co-evolutionary activity, the \emph{between} that has become, through the patient work of shared life, a home. This circle is the dynamical circle---the circle of the coupled system's trajectory through its phase space, the circle in which the holonomy of the shared life accumulates. The moment of the shared flower belongs to this circle as a relational spike: a small but real modification of the attractor landscape, a tiny increment of the holonomy that, over a lifetime of such increments, constitutes the depth of a shared life.

The outermost circle is the relational field's embeddedness in the wider social, cultural, and material world: the community, the tradition, the economic structure, the political order that either supports or undermines the conditions for intimate co-evolution. The Ubuntu framework has reminded us that the dyadic relational field is not self-sufficient but is constituted within and sustained by wider relational communities; the political-economic analysis implicit in \xref{subsec:resist} has reminded us that the conditions for relational flourishing are not given by nature but produced or destroyed by social and economic structures. The happiness of the shared flower is, ultimately, possible only in a world that provides the material, social, and cultural conditions for the cultivation of joint attention, timely sharing, and the non-possessive participation in the relational field's co-evolutionary activity.

\begin{claim}[The three circles of relational \emph{eudaimonia}]
Relational \emph{eudaimonia} has three concentric dimensions: the phenomenological (the lived moment of shared happiness, arising from the \emph{between}), the dynamical (the accumulated attractor landscape of a shared life, built through the cumulative STDP of shared contingent events), and the political-communal (the wider relational and material conditions that make the cultivation of intimate co-evolution possible). A complete account of relational happiness must attend to all three, and a complete practice of relational happiness must cultivate all three---which is to say, it is not only an intimate practice but also, irreducibly, a political and communal one.
\end{claim}

\subsection{Why This Paper Has No Synthetic Conclusion}
\label{subsec:nosynthesis}

\noindent
The reader who has reached this point may have noticed that the paper has not provided the kind of synthetic conclusion that academic papers typically provide: a paragraph or two that draws together the paper's main arguments into a unified summary, restates the paper's central thesis in its final and most complete form, and gestures toward the future work that the paper's findings make possible.

The absence of such a conclusion is not an oversight but a philosophical choice---the same choice that \pinkref{Paper~IX} made and that the series' formal律令requires \citep{paperix}. A synthetic conclusion would be a gesture toward the kind of unified, complete, fully stated truth that the GRB framework holds to be unavailable: the kind of truth that would require a metalanguage standing above all the frameworks that the paper has deployed, a view from nowhere that could survey the entire philosophical landscape and declare what the final answer is. The GRB framework holds, with Lacan and the Daoist tradition alike, that no such metalanguage exists---that the truth of the phenomenon is not to be found in any single framework's account but in the ongoing dialogue between frameworks, each of which illuminates something the others cannot see.

The happiness of the shared flower is not fully captured by the philosophical framework's account of the \emph{between}, nor by the formal framework's account of coupled oscillators and relational STDP, nor by the neuroscience of inter-brain synchrony, nor by the phenomenological description of resonance and depth, nor by the hermeneutic transformations of the classical concepts, nor by the cross-cultural encounter with \emph{ma} and Ubuntu. It is captured by all of these together, in their dialogue and their tension, in what they illuminate and in what they cannot say. The truth of relational happiness is polyphonic: it requires all the voices, and it is not reducible to any one of them.

This is the paper's final claim---not a claim about the nature of relational happiness but a claim about the nature of the philosophical inquiry into it: that such an inquiry is necessarily incomplete, that its incompleteness is not a failure but a fidelity to the complexity of the phenomenon, and that the right response to this incompleteness is not the synthesis that closes the inquiry but the continued practice---of shared attention, of timely sharing, of co-evolutionary receptivity, of non-possessive participation in the relational field's generative activity---that keeps the inquiry alive.
